\begin{frame}{추적 (1/3): 실제 경험만 있으면 $Q(S)$는 안 움직임}
\begin{center}\begin{varwidth}{0.92\textwidth}\usebeamercolor[fg]{normal text}
  \begin{itemize}
    \item 초기 $Q(S)=Q(X)=Q(T)=0$.
    \item \textbf{$Q(S)$는 안 움직임}
  \end{itemize}
  \vskip0.8em
\end{varwidth}\end{center}
\note{\begin{itemize}
\item 에이전트가 \textbf{실제로} $S\to X$(보상 0) 한 번 경험:
    \[
    Q(S) \leftarrow 0 + 0.5\,(0 + 0.9\cdot Q(X) - 0) = 0
    \]
      $Q(X)$가 아직 0이라 $Q(S)$는 \textbf{안 움직임}.
      \textbf{여기까지가 순수 Q-learning이 할 수 있는 전부}.
\item 이어서 \textbf{실제로} $X\to T$(보상 $+1$) 경험:
    \[
    Q(X) \leftarrow 0 + 0.5\,(1 + 0.9\cdot 0 - 0) = 0.5
    \]
      모델에는 $(X,a)\to(1, T)$가 추가됨.
\end{itemize}}
\end{frame}
